% Appendix C: Rigorous Definition of the Limit of a Sequence

\begin{definition}[$\epsilon$-$N$ Definition of a Limit]
Let $(a_n)$ be a sequence of real numbers. We say that $(a_n)$ \textbf{converges to a limit $L$}, written
\[
  \lim_{n \to \infty} a_n = L,
\]
if for every real number $\epsilon > 0$ there exists a positive integer $N$ such that
\[
  n > N \implies |a_n - L| < \epsilon.
\]
In formal logical notation:
\[
  \forall\, \epsilon > 0,\; \exists\, N \in \mathbb{N} \text{ such that } n > N \implies |a_n - L| < \epsilon.
\]
If no such $L$ exists, the sequence is said to \textbf{diverge}.
\end{definition}

\medskip
\noindent\textbf{Unpacking the definition.}

\begin{itemize}[leftmargin=*]
  \item \textbf{$\epsilon > 0$ is the ``margin of error.''} The challenger picks any strictly positive number,
        no matter how small. A smaller $\epsilon$ places a tighter band around $L$.

  \item \textbf{$N$ is the ``cutoff index.''} The responder must produce a concrete $N$ that works for the
        chosen $\epsilon$. Typically, smaller $\epsilon$ forces a larger $N$.

  \item \textbf{Only the \emph{tail} matters.} The condition $n > N$ means the first $N$ terms may behave
        arbitrarily; only the infinite tail of the sequence is constrained.

  \item \textbf{$|a_n - L| < \epsilon$ is the ``closeness'' condition.} Every term beyond index $N$ must
        lie strictly inside the interval $(L - \epsilon,\; L + \epsilon)$.
\end{itemize}

\medskip
\noindent\textbf{Visual interpretation (the ``$\epsilon$-tube'').}

Plot the sequence on a graph with the index $n$ on the horizontal axis and $a_n$ on the vertical axis.
Draw a horizontal dashed line at $y = L$. Choosing $\epsilon > 0$ creates a horizontal band
$\bigl(L - \epsilon,\; L + \epsilon\bigr)$ of total width $2\epsilon$ around the limit.
The definition states that for every such band there is a vertical line $x = N$ such that
\emph{every point of the sequence to the right of that line lies inside the band}.

\bigskip
\noindent\textbf{Example: Proving $\displaystyle\lim_{n \to \infty} \dfrac{1}{n} = 0$.}

\begin{proof}
Let $\epsilon > 0$ be given.

\medskip
\noindent\textit{Scratchwork.} We require
\[
  \left|\frac{1}{n} - 0\right| < \epsilon.
\]
Since $n \in \mathbb{N}$ is positive, this simplifies to $\dfrac{1}{n} < \epsilon$, i.e.\ $n > \dfrac{1}{\epsilon}$.

\medskip
\noindent\textit{Choice of $N$.} Let $N$ be any positive integer satisfying $N > \dfrac{1}{\epsilon}$.
(Such an integer exists by the Archimedean property of $\mathbb{R}$.)

\medskip
\noindent\textit{Verification.} Suppose $n > N$. Then
\[
  n > N > \frac{1}{\epsilon}
  \implies \frac{1}{n} < \epsilon
  \implies \left|\frac{1}{n} - 0\right| = \frac{1}{n} < \epsilon.
\]
Since $\epsilon > 0$ was arbitrary, the definition is satisfied and
$\displaystyle\lim_{n \to \infty} \frac{1}{n} = 0$.
\end{proof}

\medskip
\noindent\textbf{Concrete illustration.}

For $\epsilon = 0.01$ we need $N > 100$, so we may take $N = 101$.
Every term $a_n = \tfrac{1}{n}$ with $n \ge 102$ satisfies
$\tfrac{1}{102} < 0.01$, placing all those terms inside the band $(-0.01,\, 0.01)$.

\vspace{1em}
\hrule
\vspace{1em}

\noindent\textit{Note: The following advanced theorems are outside the scope of the regular HSC syllabus, but they are extremely good-to-know fundamentals for higher education and university mathematics.}

\subsection*{1. Bolzano-Weierstrass Theorem}

The Bolzano-Weierstrass Theorem is a ``guarantee of stability.'' It tells us that as long as a sequence doesn't head off to infinity, it must contain at least one ``hidden'' convergent behavior.

\begin{itemize}[leftmargin=*]
\item \textbf{The Intuition:} Imagine a sequence as a set of points trapped inside a finite box (bounded). Even if the points bounce around randomly and never settle on one spot as a whole (like $a_n = (-1)^n$), the points must eventually start ``clustering'' somewhere because they have nowhere else to go.
\item \textbf{Subsequences:} A subsequence is formed by picking infinitely many terms from the original sequence in their original order.
    \begin{itemize}
    \item \textit{Example:} Take $a_n = (-1)^n \rightarrow \{-1, 1, -1, 1, \dots\}$. The sequence itself oscillates and diverges. However, if we pick only the even terms, we get $\{1, 1, 1, \dots\}$, which converges to $1$.
    \end{itemize}
\item \textbf{The Utility:} In proofs, if you know a sequence is bounded, you can immediately ``extract'' a convergent subsequence to use for further logical steps.
\end{itemize}

\subsection*{2. Cauchy Criterion for Convergence}

The Cauchy Criterion is the ultimate ``internal'' test for convergence. 

\begin{itemize}[leftmargin=*]
\item \textbf{The Definition:} In a standard limit, we say terms get close to a \textbf{limit $L$}. In a Cauchy sequence, we say the terms get \textbf{close to each other}. 
\item \textbf{Formal Statement:} For every $\epsilon > 0$, there exists an $N$ such that for all $m, n > N$:
    \[
    |a_n - a_m| < \epsilon
    \]
\item \textbf{Why it Matters:} Usually, to prove a sequence converges, you first have to \emph{guess} what the limit $L$ is. The Cauchy Criterion removes this requirement. If you can prove that the terms of the sequence eventually ``huddle together'' and the gaps between them shrink to zero, the sequence \textbf{must} converge to some real number.
\item \textbf{The ``Completeness'' of Real Numbers:} This theorem works because the Real Number line has no ``holes.'' If the terms are huddling together, they are guaranteed to be huddling around a real number. (This is not true in the Rational numbers, where a sequence of fractions could huddle around an ``irrational hole'' like $\sqrt{2}$).
\end{itemize}
